% New theory material, incorporated verbatim into the standalone analytic source.
\part{Microscopic quantum response and moving-atom transport}
\label{part:gc-microscopic}

\section{The precise object being validated}
\label{sec:gc-contract}

A model has three distinct tests: whether its equations define an admissible
quantum process, whether its numerical solution approximates those equations,
and whether its chosen degrees of freedom, reservoirs, boundary conditions and
input parameters describe the apparatus. Passing one test does not imply the
next. The microscopic construction below supplies frequency-dependent drift,
ordered noise and photocurrent spectra for explicitly specified reduced models.
It does not identify the phenomenological mean-field gains elsewhere in this
document with a complete quantum channel. In particular, a commutator-preserving
vacuum completion of an arbitrary drift and diffusion derived from atomic jumps
are different physical models.

The local stationary four-level model, periodically driven four-level model,
finite spatial mode projection, ballistic single-atom trajectory and thermal
inflow ensemble each have a different closure assumption. Equations below make
those assumptions explicit. Internal QRT and complete-positivity tests constrain
the chosen model; they are not certificates of an experimental squeezing value.
The reported thermal ensemble is not converged, and nonlocal propagation through
the cell is not yet an established apparatus calculation.

\section{Atomic drift and diffusion from one GKSL generator}
\label{sec:gc-atomic}

Write the Hamiltonian in angular-frequency units, $h=H/\hbar$, and put rates
inside collapse operators $L_r$, with units $\mathrm{s}^{-1/2}$:
\begin{equation}
 \mathcal L\rho=-\ii[h,\rho]+\sum_r
 \left(L_r\rho L_r^\dagger-\tfrac12\{L_r^\dagger L_r,\rho\}\right).
 \label{eq:gc-gksl}
\end{equation}
Choose a complete Hermitian traceless basis $F_i$ satisfying
$\Tr F_iF_j=\delta_{ij}$, $i=1,\ldots,d^2-1$. If $\mu_i=\Tr\rho F_i$,
the adjoint equation is exactly affine on this basis:
\begin{equation}
 \dot\mu=A\mu+b,\qquad
 A_{ij}=\Tr[F_i\mathcal L(F_j)],\quad
 b_i=\Tr[F_i\mathcal L(I/d)].
 \label{eq:gc-atomic-drift}
\end{equation}
The real matrix $A$ has units $\mathrm{s}^{-1}$. This representation is exact
for one finite atom; it does not assert that the atom has a Gaussian state.
For centered operators $\delta F_i=F_i-\mu_i I$, define the ordered covariance
$C_{ij}=\Tr\rho\,\delta F_i\delta F_j$. The adjoint product rule gives
\begin{align}
 D^{(r)}_{ij}&=\Tr\rho\,[L_r^\dagger,F_i][F_j,L_r],
 &D&=\sum_rD^{(r)},\label{eq:gc-einstein}\\
 \dot C&=AC+CA^T+D. &&\label{eq:gc-covariance-identity}
\end{align}
The Hamiltonian terms cancel from the product-rule defect. Each $D^{(r)}$
is a Gram matrix: if $Q_i=[F_i,L_r]$, then
$D^{(r)}_{ij}=\Tr\rho Q_i^\dagger Q_j$. Hence $D^{(r)}\succeq0$ for
$\rho\succeq0$. Atomic $D$ has units $\mathrm{s}^{-1}$ and must not be
confused with spatial field diffusion, whose units are $\mathrm{m}^{-1}$.
For a stationary state,
\begin{equation}
 AC+CA^T+D=0,\qquad
 AK+KA^T+D-D^T=0,\quad K=C-C^T.
 \label{eq:gc-stationary-identities}
\end{equation}
The matrix $K$ is the expectation of atomic commutators and depends on the
state. It is not the constant field commutator $J$.

For a mixing stationary state, let
$R(\Omega)=(-\ii\Omega I-A)^{-1}$. With the Fourier convention of
Eq.~\eqref{eq:fourier}, the ordered spectra are
\begin{equation}
 S^>(\Omega)=RDR^\dagger,\quad
 S^<(\Omega)=RD^TR^\dagger=S^>(-\Omega)^T,\quad
 S^{\rm sym}=\tfrac12(S^>+S^<).
 \label{eq:gc-ordered-spectra}
\end{equation}
An independent full-density-matrix quantum regression calculation instead
propagates
\begin{equation}
 C_{ij}(t)=\Tr\left[\delta F_i\,e^{\mathcal L t}
                 (\delta F_j\rho_{\rm ss})\right],\quad t\geq0,
 \label{eq:gc-qrt}
\end{equation}
and adds the negative-time adjoint to its one-sided transform. Equality with
Eq.~\eqref{eq:gc-ordered-spectra} tests both source ordering and resolvent signs.
An independent probe of 30 random GKSL generators in dimensions $2,3,4$, each
with four complex jumps and five RF frequencies, found a maximum relative
ordered-spectrum discrepancy $9.99\times10^{-16}$ and maximum normalized
commutator identity residual $6.58\times10^{-16}$. These 150 comparisons
validate the finite-dimensional identities, not an alkali-vapor reservoir.
If a stationary state is nonunique or has nondecaying peripheral modes,
projecting those modes out discards elastic or preparation-dependent terms.
The stationary continuous-spectrum formula must then be supplemented by those
terms; a pseudoinverse alone is not a physical prescription.

\subsection{Reservoir admissibility and the meaning of a transit reset}
For $\rho_{\rm th}=\sum_i p_i\ketbra{i}{i}$, the operators
$L_{ij}=\sqrt{\gamma p_i}\ketbra{i}{j}$ give
\begin{equation}
 \mathcal L_{\rm reset}\rho
 =\gamma(\rho_{\rm th}\Tr\rho-\rho).
 \label{eq:gc-reset}
\end{equation}
This is a trace-preserving thermal reload channel. It is a Markov replacement
model, not a derivation of finite-flight transit times or a beam boundary
condition. In a ballistic trajectory model, applying it again merely because
the atom crosses numerical segments would reset the same atom repeatedly.

Trace preservation, Hermiticity preservation and conditional complete
positivity of the generator must be checked separately. In the Choi test,
project the generator Choi matrix onto the space orthogonal to the maximally
entangled vector and require positive semidefiniteness there. Positivity of a
particular stationary density matrix does not imply this property. A
coherence-only transit-damping variant has a projected-Choi eigenvalue about
$-0.00607$ in its normalized generator audit and a normalized Choi eigenvalue
about $-3.08\times10^{-4}$ after $1\,\mathrm{ns}$. This counterexample excludes
that variant as a general GKSL reservoir even when its computed steady state
looks admissible. Passing a radiative or selected thermal-reset model does not
validate every phenomenological collision law.

\section{Reciprocal SI normalization and hyperfine reduction}
\label{sec:gc-normalization}

For a Gaussian beam with $1/\e^2$ intensity radius $w$,
\begin{equation}
 I_0=\frac{2P}{\pi w^2},\quad
 E_0=\sqrt{\frac{2I_0}{\epsilon_0c}},\quad
 \Omega=\frac{dE_0}{\hbar},\quad h_{eg}=\frac{\Omega}{2}.
 \label{eq:gc-si-drive}
\end{equation}
The same signed dipole must enter the drive, emitted polarization and
radiative normalization. In the scalar D1 convention used here,
\begin{equation}
 d_\Gamma^2=\frac{3\pi\epsilon_0\hbar c^3\Gamma}{\omega_{D1}^3}.
 \label{eq:gc-dipole-linewidth}
\end{equation}
It gives $d_\Gamma=2.5367896004\times10^{-29}\,\mathrm{C\,m}$ for the
declared constants. A stored value $2.5377\times10^{-29}\,\mathrm{C\,m}$
implies a linewidth $0.0717886\%$ larger. There is no additional universal
$\sqrt2$ multiplier: a change of reduced-matrix-element convention requires
changing the associated angular factors too. A diagnostic pump/weak-field
comparison gives implicit dipoles $1.4646827595\times10^{-29}$ and
$7.3257088906\times10^{-30}\,\mathrm{C\,m}$, a ratio $1.9993734141$.
One cannot repair such a mismatch by fitting output gain alone.

Let $S_{FF'}$ be hyperfine strengths normalized within a ground manifold.
For an unpolarized manifold with population $p_F$, weak absorption in a fixed
linear polarization has the population-weighted factor
\begin{equation}
 \frac{3C_F^2p_F}{2F+1}=\frac{p_FS_{FF'}}{3},
 \qquad 3C_F^2=\frac{(2F+1)S_{FF'}}{3}.
 \label{eq:gc-manifold-strength}
\end{equation}
For the four D1 branches $(2,2),(2,3),(3,2),(3,3)$, the required per-manifold
factors $S_{FF'}/3$ are $(2,7,5,4)/27$. The alternative common-factor
prescription $3C_F^2/12$ gives $(10,35,35,28)/324$, with ratios $5/12$,
$5/12$, $7/12$, $7/12$ to those factors. Thus no single fitted scalar can
replace both manifold denominators $5$ and $7$ when $p_2,p_3$ evolve
independently. For a thermal distribution $p_F=(2F+1)/12$, the exact weak
absorption factor is $3C_F^2/12$ \emph{after} including population once.
Multiplying it by $p_F$ again is incorrect.

A four-state surrogate may use $d_{FF'}=d_\Gamma\sqrt{S_{FF'}/3}$ with
trace-one manifold populations. It reproduces this weak absorption average,
but not the signs, coherences, polarization selection or optical pumping of
the full Zeeman Hamiltonian under a strong drive. The residual $\ell_s=0.74$
in the separate phenomenological model is therefore an assumption, not a
microscopic line-strength determination.

The optical carrier must also refer to the same transition as the detuning.
For the specified ground and excited splittings, the transition anchor differs
from the fine-structure centroid by
$(7\omega_{\rm hf,g}+5\omega_{\rm hf,e})/12$, or
$1.921502256083\GHz$ in ordinary frequency. Rounding the data tables can
leave a $1\kHz$ splitting discrepancy. In one weak-gain fixture a carrier-only
change moves $G_p$ from $1.0517029257$ to $1.0517034724$; its smallness in that
fixture does not justify mixing carrier conventions in a derivation.

\section{Eliminating the atom without guessing field noise}
\label{sec:gc-field}

Let $n$ denote number density, $A_\perp$ the collected area and
$\lambda=nA_\perp$ the line density. Write the field amplitude as
$E_j=Q_ja_j$, $Q_j=\sqrt{2\hbar\omega_j/(\epsilon_0cA_\perp)}$ and define
$g_j=d_jQ_j/(2\hbar)$. Then $g_j$ has units $\mathrm{s}^{-1/2}$.
For the Nambu vector $b=(a_p,a_c^\dagger)^T$, choose atomic operators $O_j$
consistently with its creation/annihilation ordering. A reciprocal construction
has, with signs fixed by that choice,
\begin{equation}
 W_{ji}=\Tr(O_jF_i),\quad
 B_{ij}=-\ii g_j\Tr\{F_i[O_j^\dagger,\rho_{\rm ss}]\},\quad
 C_0=-\ii J\operatorname{diag}(g_j)W.
 \label{eq:gc-reciprocal-couplings}
\end{equation}
In a stationary local medium the linearized equations and the slice-averaged
atomic source normalization are
\begin{equation}
 \partial_t\delta F=A\delta F+Bb+\xi,\quad
 \partial_zb=\lambda C_0\delta F+K_{\rm geom}b,\quad
 \langle\xi(z,t)\xi(z',t')^T\rangle
 =\frac{D}{\lambda}\delta(z-z')\delta(t-t').
 \label{eq:gc-local-equations}
\end{equation}
Eliminating the atom at the appropriate rotating-frame frequency gives
\begin{align}
 M(\Omega)&=\lambda C_0R(\Omega)B+K_{\rm geom},\label{eq:gc-field-drift}\\
 D_f^>(\Omega)&=\lambda C_0R(\Omega)DR(\Omega)^\dagger C_0^\dagger,
 \qquad
 D_f^<(\Omega)=\lambda C_0R(\Omega)D^TR(\Omega)^\dagger C_0^\dagger.
 \label{eq:gc-field-diffusion}
\end{align}
The conjugate-frequency sectors must be aligned before assembling the physical
sideband matrix. Both drift and field diffusion are linear in density. For
one common area, $\lambda g_ig_j$ is independent of that area; the coherent
seed amplitude still scales as $\sqrt{P/\hbar\omega}$. Independent velocity
classes contribute weighted drift and weighted covariance, not the square of
an averaged noise amplitude. Every bath must appear once; radiation already
represented by an explicitly retained field cannot also be counted as an
independent unobserved reservoir without a consistent mode partition.

The local approximation assumes short-range independent atomic fluctuations
at different $z$. It is unsuitable for the same ballistic atom observed in
successive illuminated regions. Sections~\ref{sec:gc-transport}--\ref{sec:gc-ensemble}
replace its spatial delta function by an explicit two-position covariance.

\section{Physical sidebands, complete positivity and finite RF modes}
\label{sec:gc-sidebands}

At a strictly positive analysis frequency the physical annihilation modes are
$(a_{p,+},a_{c,-},a_{p,-},a_{c,+})$, with eight real quadratures. A complex
$2\times2$ Nambu covariance is only one sector and is not by itself a
two-oscillator covariance to which an arbitrary real symplectic test may be
applied. Given a sector transfer $T$, its Bogoliubov matrices are
$U=\operatorname{diag}(T_{00},T_{11}^*)$ and
$V_{01}=T_{01}$, $V_{10}=T_{10}^*$, with the other entries zero. In an
$(X_1,X_2,P_1,P_2)$ ordering its real transfer is
\begin{equation}
 X=\begin{pmatrix}
 \Rea(U+V)&-\Ima(U-V)\\
 \Ima(U+V)&\Rea(U-V)
 \end{pmatrix}.
 \label{eq:gc-bogoliubov-real}
\end{equation}
Permute it into the declared quadrature ordering and include both RF sectors.
If $H=(N^>+N^<)/2$ is the integrated sector noise, its anomalous off-diagonal
entry $m=H_{01}$ contributes the real cross block
$\left(\begin{smallmatrix}\Rea m&\Ima m\\\Ima m&-\Rea m\end{smallmatrix}\right)$;
the diagonal blocks are $H_{00}I_2,H_{11}I_2$.

For canonical quadratures with vacuum $I/2$ and
$J_R=\bigoplus_j\left(\begin{smallmatrix}0&1\\-1&0\end{smallmatrix}\right)$,
the Gaussian channel and state conditions are respectively
\begin{equation}
 Y+\frac{\ii}{2}(J_R-XJ_RX^T)\succeq0,
 \qquad V+\frac{\ii}{2}J_R\succeq0.
 \label{eq:gc-cp}
\end{equation}
These are distinct from merely requiring $Y\succeq0$, and from the Nambu
commutator identity. They use the canonical normalization of
Ref.~\cite{Weedbrook2012}.

A finite measured RF mode is
$a_h=\int\dd\Omega\,h(\Omega)^*a(\Omega)/(2\pi)$ with
$\int\dd\Omega\,|h|^2/(2\pi)=1$. A top-hat interval of width $B$ in Hz
therefore has $|h|=1/\sqrt B$. Its time profile has sinc tails and is not a
compact time gate. Tracing orthogonal vacuum modes after frequency-dependent
transfer adds more than the average $Y$:
\begin{equation}
 \overline X=\sum_kp_kX_k,\qquad
 Y_{\rm eff}=\sum_kp_kY_k+\tfrac12\sum_kp_k
             (X_k-\overline X)(X_k-\overline X)^T,
 \quad\sum_kp_k=1.
 \label{eq:gc-filtered-channel}
\end{equation}
For equally weighted rotations by $\pm\theta$, the retained mode has
$\overline X=\cos\theta I$, $Y_{\rm eff}=\sin^2\theta I/2$. Omitting the
second term would turn a passive physical process into a CP violation.
Equation~\eqref{eq:gc-filtered-channel} assumes the discarded input modes are
vacuum and initially independent of the retained mode; a nonvacuum multimode
input requires its full covariance.

\section{Finite coherent seeds and periodic quantum regression}
\label{sec:gc-periodic}

A finite coherent seed makes the reference atom periodic. In a declared signed
frequency convention write $h(t)=h_0+Ve^{-\ii\nu t}+V^\dagger e^{\ii\nu t}$,
with $\nu=-\omega_{\rm hf}+\delta$ for the chosen Raman branch, and expand
$\rho(t)=\sum_q\rho_qe^{-\ii q\nu t}$. The diffusion in
Eq.~\eqref{eq:gc-einstein} is evaluated on $\rho(t)$, not just $\rho_0$.
The exact time-dependent identity is still
$\dot C=A(t)C+CA(t)^T+D(t)$.
For harmonic indices $h,k$, the lifted matrices are
\begin{equation}
 \mathbb A_{hk}=A_{h-k}+\ii h\nu\delta_{hk}I,\quad
 \mathbb D_{hk}=D_{h-k},\quad
 \mathbb R(\Omega)=(-\ii\Omega I-\mathbb A)^{-1}.
 \label{eq:gc-lifted-generator}
\end{equation}
The ordered spectrum is $\mathbb R\mathbb D\mathbb R^\dagger$. To form the
lesser source, transpose the \emph{atomic operator indices} in each Fourier
coefficient; transposing the entire harmonic matrix changes the harmonic
ordering and is generally wrong. A shift of time origin rotates cross-harmonic
blocks covariantly rather than deleting them.

The density-state cutoff and response cutoff are separate approximation
parameters. Check interior physical blocks under enlargement of both; positivity
alone cannot certify harmonic convergence. The exact one-period propagator is
CPTP and a mixing periodic solution has spectral radius below one on the
traceless subspace. A finite harmonic matrix need not itself be a GKSL
generator. An independent density-matrix period reference uses
\begin{equation}
 Q(\Omega)=\int_0^T\dd s\,e^{\ii\Omega s}\mathcal U(t_0+s,t_0),\quad
 \int_0^\infty\dd t\,e^{\ii\Omega t}\mathcal U(t_0+t,t_0)
 =Q(\Omega)[I-e^{\ii\Omega T}\mathcal P(t_0)]^{-1},
 \label{eq:gc-period-qrt}
\end{equation}
with trace bordering/projection on the decaying subspace. This reference does
not obtain its noise by copying the Floquet diffusion matrix.

For finite propagating mean fields $\beta$, the mean equation is nonlinear,
$\beta'=F(\beta,\beta^*)$. Its zero-frequency fluctuation transfer is the
tangent map of this equation. In general the mean output is not obtained by
applying that tangent transfer to the input mean. Holding the pump prescribed
also remains different from solving three-field pump depletion.

For the reciprocal four-state RMS-dipole fixture with
$n=10^{18}\,\mathrm{m}^{-3}$, $A_\perp=1.2\times10^{-7}\,\mathrm{m}^2$,
$L=12.5\,\mathrm{mm}$, pump $0.6\,\mathrm W$, $w=530\,\mu\mathrm m$,
$\Delta/(2\pi)=0.9\GHz$, $\delta/(2\pi)=-8\MHz$,
reset rate $2\pi\,100\kHz$, and illustrative efficiency $0.85$, the
stationary reference gives $G_p=1.11065587$, $G_c=0.11356221$ and a detected
ratio $-0.7302943383\dB$ at $1\MHz$. These inputs specify a mathematical
fixture, not measured density, collisions or detector calibration.
An $8\,\mu\mathrm W$ periodic-seed calculation changes this ratio by about
$10^{-4}\dB$, while deleting cross-sector correlations changes a $0.1\MHz$
example by $0.039\dB$. Individual noise components can differ by about
$40\%$ despite a small final ratio change. These observations do not give a
global finite-seed error bound.

Gaussian photocounting adds a separate fourth-moment closure to the exact
atomic second-moment QRT. Beyond the bright-seed linearization, the spontaneous
mean and quadratic fluctuation intensity enter both PSD and SQL. A declared
finite-band calculation at $8\,\mu\mathrm W$ changes the bright-seed result
by $1.033\times10^{-5}\dB$, but increasing the optical collection bandwidth
from $512$ to $1024\MHz$ increases spontaneous collected flux by $5.016\%$.
Thus no bound on an unfiltered optical tail follows from the small finite-band
photocurrent correction.

\section{Two moving phases and spatial optical modes}
\label{sec:gc-kinetic}

Define laboratory wavevectors $\mathbf k_0,\mathbf k_p,\mathbf k_c$ and
$\mathbf q=\mathbf k_p-\mathbf k_0$,
$\mathbf Q=\mathbf k_p+\mathbf k_c-2\mathbf k_0$. Two phase coordinates are
\begin{equation}
 \theta_1=\nu t-\mathbf q\cdot\mathbf r,\qquad
 \theta_2=-\mathbf Q\cdot\mathbf r,\qquad
 \dot{\boldsymbol\theta}
 =(\nu-\mathbf q\cdot\mathbf v,-\mathbf Q\cdot\mathbf v).
 \label{eq:gc-moving-phases}
\end{equation}
For couplings $V_p,V_c$, a consistent Hamiltonian is
$h=h_0+V_pe^{-\ii\theta_1}+V_ce^{-\ii(-\theta_1+\theta_2)}+\mathrm{h.c.}$.
Expanding on the integer two-phase lattice yields
\begin{equation}
 0=\sum_{\mathbf b}\mathcal L_{\mathbf a-\mathbf b}\rho_{\mathbf b}
    +\ii(\mathbf a\cdot\dot{\boldsymbol\theta})\rho_{\mathbf a}.
 \label{eq:gc-kinetic-lattice}
\end{equation}
At exact $\mathbf Q=0$, harmonics that represent the same physical phase
must be coherently identified: combine $V_p+V_c^\dagger$ before solving.
They are not independent photon ports. If $\mathbf Q\cdot\mathbf v=0$
but $\mathbf Q\ne0$, the second phase is a static spatial label and cannot
simply be discarded.

This distinction has an observable consistency test. In a fixture with
second-phase frequency $-1.123879\times10^6\,\mathrm{s}^{-1}$, freezing the
moving grating omits the term $\dot\theta_2\partial_{\theta_2}K$ from the
commutator identity. Diffusion remains PSD, but the normalized commutator
defect is $2.96091344\times10^{-5}$. Increasing the lattice cutoff changes
that residual by only $4.14\times10^{-15}$: the omission is physical, not a
resolution error.

For spatial profiles $u_j(\mathbf r_\perp)$ normalized by
$\int |u_j|^2\dd^2r_\perp=1$, the mode-projected response integrates the
local polarization with both drive and collection profiles. Cross-mode noise
comes from shared atoms and must retain those cross correlations. The limit
$\mathbf Q\to0$ is taken at fixed finite aperture, not by replacing its
Fourier transform with an infinite-area delta function. A two-mode truncation
does not establish convergence in the number of transverse modes.

\begin{table}[ht]
\centering\small
\caption{Conditional stationary spatial fixture: a $400\times300\,\mu\mathrm m$
rectangular aperture, prescribed pump, and the reciprocal model parameters of
Sec.~\ref{sec:gc-periodic}. Noise ratios refer to $1\MHz$.}
\label{tab:gc-spatial}
\begin{tabular}{rrrr}
\toprule
Angles $(\theta_p,\theta_c)$ in mrad&$G_p$&$G_c$&$S_-\ (\dB)$\\
\midrule
$(0,0)$&1.1106391761&0.1135452814&$-0.7301956922$\\
$(5,-5)$&1.0620337919&0.0650877602&$-0.4328093597$\\
$(5,-4)$&1.0277727399&0.0306514955&$-0.2040641891$\\
\bottomrule
\end{tabular}
\end{table}
For the asymmetric-angle case, a transverse quadrature change $N_x=4\to8$
moves the ratio by $3.5784\times10^{-6}\dB$ and fails a declared
$10^{-6}\dB$ tolerance; $8\to12$ changes it by $3.93\times10^{-13}\dB$.
Independent local QRT comparisons give relative noise and drift residuals
$3.86\times10^{-12}$ and $1.58\times10^{-13}$; the DC mean-field tangent
test gives $5.77\times10^{-10}$ and the global commutator test
$1.38\times10^{-11}$. Only the RF points $0.1,1,4\MHz$ are sampled. Neither
these three points nor transverse quadrature convergence establishes a
bandwidth or a full spatial mode basis.

\section{Ballistic finite-age atoms and nonlocal covariance}
\label{sec:gc-transport}

An atom enters at $\mathbf r_{\rm in}$ with velocity $\mathbf v$ and
prepared state $\rho_{\rm in}$. Its age $a$ labels
$\mathbf r(a)=\mathbf r_{\rm in}+\mathbf va$ and
$\dot\rho(a)=\mathcal L(a)\rho(a)$ until its first exit at $a=\tau$.
Let $\Phi(a,b)$ be the atomic fluctuation propagator. Then
\begin{align}
 C(a,b)={}&\Phi(a,0)C_{\rm in}\Phi(b,0)^T\nonumber\\
 &+\sum_r\int_0^{\min(a,b)}\dd s\,
      \Phi(a,s)D^{(r)}(s)\Phi(b,s)^T.
 \label{eq:gc-transport-covariance}
\end{align}
The first term is initial internal quantum noise; the second is noise from
jumps during the flight. Both contain correlations between different ages.
The field response is also retarded, containing
$\Theta(a-b)C_{\rm out}(a)\Phi(a,b)B_{\rm drive}(b)$. A diagonal-in-age
noise replacement loses shared-atom correlations even when each retained
diagonal block is PSD.

One computes the integrated atomic output
\begin{equation}
 Y_j=\int_0^\tau\dd a\,e^{\ii\omega_j a}O_j(a)
 \label{eq:gc-integrated-output}
\end{equation}
by augmenting the fluctuations with $Y/\tau$. Its drift has blocks
$\left(\begin{smallmatrix}A&0\\e^{\ii\omega a}C_{\rm out}/\tau&0\end{smallmatrix}\right)$
and its direct diffusion has $D$ only in the atomic block. The integrated
covariance is the final output block multiplied by $\tau^2$; it includes all
off-diagonal age correlations without storing an age-by-age kernel.

An independent adjoint formulation defines operators
\begin{equation}
 \mathcal A_j(t)=\int_t^\tau\dd s\,
  \mathcal U(s,t)^\dagger[O_j]e^{\ii\omega_js},\quad
 \dot{\mathcal A}_j=-\mathcal L(t)^\dagger\mathcal A_j-e^{\ii\omega_jt}O_j,
 \quad\mathcal A_j(\tau)=0.
 \label{eq:gc-adjoint-output}
\end{equation}
For $\mu_j=\Tr\rho_{\rm in}\mathcal A_j(0)$, its source-resolved greater
covariance is
\begin{align}
 C_{{\rm in},jk}^>&=\Tr[\rho_{\rm in}\mathcal A_j(0)\mathcal A_k(0)^\dagger]-\mu_j\mu_k^*,
 \label{eq:gc-adjoint-boundary}\\
 C_{r,jk}^>&=\int_0^\tau\dd t\,
 \Tr\rho(t)[L_r^\dagger,\mathcal A_j(t)]
                   [\mathcal A_k(t)^\dagger,L_r].
 \label{eq:gc-adjoint-jumps}
\end{align}
The lesser covariance reverses the operator ordering. Linear response to an
independent perturbation $V_ke^{-\ii\nu_kt}$ follows from
\begin{equation}
 \mathcal R_{jk}=-\ii\int_0^\tau\dd t\,
 e^{-\ii\nu_kt}\Tr\{\rho(t)[\mathcal A_j(t),V_k]\}.
 \label{eq:gc-adjoint-response}
\end{equation}
This reference constructs source covariances directly from commutators and
the full density matrix, independently of the forward Einstein-matrix
assembly. A separate two-time density-matrix QRT integral tests the complete
integrated covariance. In a boundary-noise deletion counterexample all
retained noise matrices stay PSD, but the commutator defect is $0.20196$ and
the output spectrum changes by $14.59\%$. Likewise, a diagonal-only kernel
sum decreases toward zero under refinement whereas the full two-time
trapezoidal integral converges to a finite covariance.

\section{Thermal inflow, random arrivals and the unresolved cell closure}
\label{sec:gc-ensemble}

For independent Poisson arrivals with rate $\mathcal J$, a single-path
integrated mark with connected covariance $C$ and mean $\mu$ has spectrum
\begin{equation}
 S_{\rm stream}=\mathcal J(C+\mu\mu^\dagger),\qquad
 \overline N=\mathcal J\tau.
 \label{eq:gc-arrival-noise}
\end{equation}
The outer product is arrival-number noise, distinct from the internal atomic
boundary covariance. For $O=I$ the internal noise vanishes, but random entry
times still give $\mathcal J|\int_0^\tau e^{\ii\omega a}\dd a|^2$.
With one random common entry phase, charges $(1,-1,-1,1)$ retain only
equal-charge cross terms. Average $\mu\mu^\dagger$ before subtracting or
combining means; a zero phase-averaged $\mu$ does not imply zero number noise.

For a Maxwell velocity density $f_M$ normalized in three dimensions and an
inward boundary normal, the incoming flux measure is
\begin{equation}
 \dd\mathcal J=n f_M(\mathbf v)|\mathbf v\cdot\mathbf n|
   \dd^2r\,\dd^3v,\qquad
 \int\tau(\mathbf r,\mathbf v)\dd\mathcal J=nV.
 \label{eq:gc-thermal-flux}
\end{equation}
The first-exit lifetime $\tau$ depends on the boundary point and velocity.
The occupancy identity checks geometric flux weights, not the oscillatory
response/noise integral. Quadrature weights multiply path covariances once;
squaring those weights or inserting density a second time is incorrect.

The declared thermal fixture is a square column of side
$\sqrt{1.2\times10^{-7}}=346.410\,\mu\mathrm m$, length $12.5\,\mathrm{mm}$,
temperature $373\,\mathrm K$ and separately specified density
$10^{18}\,\mathrm{m}^{-3}$. It injects an unpolarized internal state
$\operatorname{diag}(5/12,7/12,0,0)$ at the box boundary, with no upstream
pumping history, re-entry or velocity-changing collisions. This is a
mathematical boundary-value problem, not automatically the physical boundary
of a vapor cell. With a $530\,\mu\mathrm m$ pump waist, its boundary pump
intensity is $0.652$--$0.808$ of peak. A frozen zero-velocity atom at a side
midpoint, exposed to the declared $0.6\,\mathrm W$ pump and $0.9\GHz$
detuning, changes its $F=3$ population from $0.583333$ to $0.645138$,
$0.869748$, and $0.942976$ after $0.1,1,5\,\mu\mathrm s$. The corresponding
Frobenius distances from the injected density matrix are $0.20293,0.42950,
0.52642$. This is a boundary-sensitivity example, not an estimate of the
experimental modeling error; it shows why entry-state preparation requires
physical justification.

In a transported ensemble the optical equation contains a spatially nonlocal
susceptibility kernel, and the output noise contains a double spatial
integral:
\begin{align}
 \partial_z b(z)&=K_{\rm geom}b(z)+
   \int\dd z'\,\mathcal R(z,z';\Omega)b(z')+\eta(z),\nonumber\\
 W_{\rm out}&=\int\dd z\,\dd z'\,
 G(L,z)\mathcal C(z,z';\Omega)G(L,z')^\dagger.
 \label{eq:gc-nonlocal-maxwell}
\end{align}
An integrated finite-age response is not, by dimensional relabeling, a local
matrix $M(\Omega)$. Nor may correlated trajectory segments be composed as
independent local Gaussian channels. A solved, converged nonlocal optical
boundary problem is required before a thermal transport spectrum becomes a
cell-output squeezing prediction.

\section{Exact constant-step convolution and controlled trajectory numerics}
\label{sec:gc-numerics}

For a constant generator step, density evolution and source-resolved covariance
can be integrated simultaneously. The scalar convolution
\begin{equation}
 F(x,y;h)=\int_0^h\dd u\,e^{x(h-u)}e^{yu}
 =\begin{cases}(e^{yh}-e^{xh})/(y-x),&y\ne x,\\
 he^{xh},&y=x,
 \end{cases}
 \label{eq:gc-divided-exponential}
\end{equation}
is evaluated with a stable divided-exponential expression near coincidence.
If the density generator has eigenvalues $\lambda_k$, the augmented
fluctuation drift has eigenvalues $\beta_a$, and each source $D_r(\rho)$
is linear in density, the particular covariance solution in that eigenbasis
has entries
\begin{equation}
 (\widetilde C_r(h))_{ab}
 =e^{(\beta_a+\beta_b^*)h}(\widetilde C_r(0))_{ab}
 +\sum_k\widetilde D_{r,abk}\,q_k
      F(\beta_a+\beta_b^*,\lambda_k;h).
 \label{eq:gc-exact-step}
\end{equation}
The eigenoperators of the density generator can be complex and nonpositive.
Hermitizing, clipping their eigenvalues, or normalizing their traces destroys
the linear expansion. A robust implementation checks eigenbasis condition
numbers (a declared guard $10^6$), a roundoff amplification estimate
$\epsilon\kappa(S)^2\kappa(R)\leq10^{-8}$, and eigen residuals
$\leq10^{-12}$; otherwise it uses a block exponential or another stable
coordinate-free solve. Exact integration of each constant step does not make
a changing pump envelope exact.

For an envelope $f(t)$ multiplying one Hamiltonian term, a fourth-order
commutator-free composition uses nodes $c_{1,2}=1/2\mp\sqrt3/6$ and
$a_{1,2}=(3\mp2\sqrt3)/12$:
\begin{align}
 f_A&=2(a_2f(c_1h)+a_1f(c_2h)),\quad
 f_B=2(a_1f(c_1h)+a_2f(c_2h)),\nonumber\\
 U_{\rm CF4}(h)&=\exp[\tfrac h2(G_0+f_BG_1)]
                 \exp[\tfrac h2(G_0+f_AG_1)].
 \label{eq:gc-cf4}
\end{align}
Here the entire affine density/covariance lift follows the same composition.
Although $a_1<0$, each duration is positive and each substep uses a Hermitian
Hamiltonian with the same positive-rate jumps, so that each atomic substep is
GKSL. Clipping a combined envelope changes the method and its order.
General background on such propagation is given in
Ref.~\cite{Alvermann2012}.

A declared path gate requires three primary resolutions with two successive
relative changes below $10^{-3}$; agreement with a separately refined
reference below $5\times10^{-6}$; and reference self-refinement below
$2\times10^{-6}$. Gates apply to each RF matrix and each physical noise
source, with a denominator floor $128\epsilon$ times a declared SI scale.
That floor stabilizes a relative residual; it adds no physical noise.
A demanding $2\,\mu\mathrm s$ trajectory changes by $9.4185\times10^{-4}$
from 4096 to 8192 CF4 steps, then by $1.6138\times10^{-4}$ at 16384 steps,
and differs from its independent reference by $8.6908\times10^{-7}$.
For the separately sampled seed-811 path set, maximum successive changes are
$5.347\times10^{-6}$ and $2.135\times10^{-7}$, the forward/adjoint
discrepancy is $6.08\times10^{-8}$, and reference refinement is
$1.385\times10^{-7}$. These are path-solver results, not ensemble convergence.

\subsection{A converged path is not a converged thermal ensemble}
The current thermal sample contains three independent scrambles, labelled
11, 211 and 811, of the same coarse $p=2$ quadrature: 72 trajectories and 360
primary/reference solves in total. Directional relative comparisons use the
second sample as denominator, so both directions are tested. The range across
the six ordered comparisons is:
\begin{table}[ht]
\centering\small
\caption{Thermal ensemble differences for the declared coarse inflow problem.
Every ordered seed-pair comparison fails the $5\%$ all-observable gate.}
\label{tab:gc-thermal-convergence}
\begin{tabular}{lrr}
\toprule
Quantity&Minimum difference (\%)&Maximum difference (\%)\\
\midrule
Greater ordered covariance&12.8553&19.9038\\
Lesser ordered covariance&12.8928&19.9644\\
Greater covariance, by source&25.5642&43.6110\\
Lesser covariance, by source&22.1673&43.3337\\
Arrival-number contribution&5.6997&9.0633\\
Linear response&25.7841&34.3654\\
\bottomrule
\end{tabular}
\end{table}
The occupancy estimates divided by $nV$ are $0.8302268543$, $1.1102815979$
and $0.8561468387$. Path accuracy can be good while the ensemble is
under-resolved; agreement of one integral cannot validate the others. There
is no accepted refinement sequence for all response, ordered covariance,
source-resolved covariance and number-noise matrices. Consequently neither
the thermal covariance nor a derived SQL-normalized squeezing spectrum is
certified by these samples.

\section{Readout precision and uncertainty conditional on the model}
\label{sec:gc-readout}

For detected currents $I_p,I_c$, complex electronic transfer functions $H_p,H_c$
and subtraction weight $w$, the one-sided shot-noise PSD is
\begin{equation}
 S_{\rm SQL}^{(1)}(f)=2e\bigl(|H_p(f)|^2I_p+w^2|H_c(f)|^2I_c\bigr).
 \label{eq:gc-sql}
\end{equation}
The quantum photocurrent quadratic form uses the same mean currents,
sideband ordering and electronic phases. Electronics noise and its
cross-correlations are added with their measured signs. With
$\langle x^2\rangle=\int S_\Omega(\Omega)\dd\Omega/(2\pi)
=\int S_f(f)\dd f$, the two-sided convention is
$S_f(f)=S_\Omega(2\pi f)$: no extra $2\pi$ appears. For a real signal the
positive-frequency one-sided PSD is twice that two-sided value. Numerator
and SQL must use the same convention.

Average linear PSD and SQL through the actual RF resolution filter before
forming their ratio, and convert to $10\log_{10}$ only afterward. Averaging
dB values is a different observable. Finite bright-seed, Gaussian fourth-moment
and optical collection-band assumptions are stated in
Sec.~\ref{sec:gc-periodic}; they are not consequences of the detector SQL
identity.

For independently characterized parameters $p$ with covariance $\Sigma_p$,
linear propagation gives $\Sigma_y=J_p\Sigma_pJ_p^T$, where
$(J_p)_{ij}=\partial y_i/\partial p_j$. Joint sampling is preferable when
responses are nonlinear, but must preserve parameter correlations,
nonnegative rates and physically allowed density matrices. A parameter fixed
by assumption is not independently measured merely because it was not fitted
to the final squeezing trace. Uncertainty intervals conditioned on an
unconverged thermal grid or unjustified boundary state do not cover those
model errors. Measured gain, bandwidth, mode overlap, phase response and
detector calibration must be held out from any fit used to claim predictive
agreement.
